\section{既存の方法や理論}
\if0
\UTF{2022} 似たような（あるいは同じ）目的のシステムと目指したシステムの原理
的・技術的な違い
\UTF{2022} 開発したシステムの原理の採用理由や長所・短所など
\UTF{2022} 盛り込んだ技術的工夫の新規性や独自性
\fi
複数機器の同期制御に関する既存技術として，MIDI クロック同期，Ableton Link，DMX512 が挙げられる．MIDI クロック同期は，1台のマスターデバイスがクロック信号を複数のスレーブ機器へ送信し，ドラムマシンやシンセサイザなどのテンポを一致させる方式である．Ableton Link は，ローカルネットワーク上の複数の PC やスマートフォンで動作する音楽アプリケーション間のテンポや位相を同期させるプロトコルである．また，DMX512 は，照明制御を目的とした1対多の一斉送信プロトコルであり，1台のコントローラから多数の機器を制御できる．これらの技術は，複数機器を同期させる代表的な方法であり，本システムと比較する対象となる．

\begin{table}[H]
\centering
\caption{既存技術の概要と課題}
\label{tab:1-1}
\begin{tabularx}{\linewidth}{|>{\centering\arraybackslash}p{3.0cm}|X|X|}
\hline
\textbf{名称} & \textbf{概要} & \textbf{課題} \\
\hline
MIDIクロック同期
&
マスター1台から複数のスレーブ機器へクロック信号を送り，テンポを同期させるプロトコルである．
&
デイジーチェーン接続では台数の増加に伴って遅延が蓄積する可能性がある．また，専用の MIDI インターフェースが必要となる．
\\
\hline
Ableton Link
&
ローカルネットワークを介して，複数のアプリケーション間でテンポや位相を同期させるプロトコルである．
&
Wi-Fi などのネットワーク環境が必須であり，Arduino などのマイコン単体では動作させにくい．
\\
\hline
DMX512
&
照明制御用の1対多一斉送信プロトコルであり，マスター1台から最大512台の機器を制御できる．
&
専用の RS-485 トランシーバ IC や専用ケーブルが必要となり，小規模な演奏制御システムでは構成が大きくなりやすい．
\\
\hline
\end{tabularx}
\end{table}

MIDI クロック同期は，音楽機器のテンポ同期を目的とした技術であり，本システムと目的が近い．しかし，MIDI 機器を前提とするため，各デバイスに MIDI インターフェースが必要となる．また，デイジーチェーン接続を用いる場合，接続台数が増えることで信号伝達の遅延が蓄積する可能性がある．本プロジェクトでは，Arduino を複数台用いた小規模な構成を対象としているため，専用の MIDI 機器を追加する方法は，コストや実装の容易さの面で適していない．

Ableton Link は，PC やスマートフォン上の音楽アプリケーションをネットワーク越しに同期できるため，高機能な音楽同期手法である．しかし，Wi-Fi などのローカルネットワーク環境を前提としており，Arduino のようなマイコン単体で直接扱うには適していない．また，本システムでは，音楽アプリケーション同士の同期ではなく，Master Arduino から複数の Slave Arduino へ設定情報を送信し，各 Slave が担当楽器の演奏制御を行うことを目的としている．そのため，Ableton Link のようなアプリケーション間同期よりも，マイコン間で扱いやすい通信方式が必要である．

DMX512 は，1台のコントローラから多数の機器へ制御情報を送信できる点で，本システムの1対多制御と類似している．一方で，DMX512 は主に照明制御を目的とした規格であり，RS-485 トランシーバ IC や専用ケーブルを必要とする．本プロジェクトでは，教育現場やハッカソンのような限られた機材環境でも再現しやすい構成を重視しているため，専用部品を追加する方式は構成が大きくなりやすい．

これらの既存技術に対し，本システムでは Arduino に標準的に搭載されている I2C 通信を採用する．I2C 通信は，SDA と SCL の2本の信号線を共有するバス型の通信方式であり，1台の Master から複数の Slave へ情報を送信できる．本システムでは，Master が演奏順番，オクターブ，BPM などの設定情報を送信し，各 Slave が受信した情報に基づいて担当楽器の演奏を制御する．I2C を用いることで，専用の MIDI インターフェースや RS-485 トランシーバ IC を用いずに，Arduino 本体，プルアップ抵抗，ジャンパワイヤを中心とした低コストな同期基盤を構成できる．

I2C 採用の長所は，少ない配線で複数台の Arduino を接続できる点，Arduino 標準のライブラリで扱いやすい点，Master と Slave の役割を明確に分けられる点である．これにより，Master は設定情報の送信と全体制御を担当し，各 Slave は担当楽器の楽譜データや発音処理を担当するという処理分担が可能となる．一方で，I2C は本来，同一基板内や短距離通信での利用を想定した通信方式であるため，長距離配線やノイズの多い環境では通信が不安定になる可能性がある．したがって，本システムでは，小規模な卓上システムとして利用することを前提とし，配線長や接続状態に注意して設計する必要がある．

本システムの技術的工夫は，Arduino 標準機能である I2C 通信を用いて，複数の Slave に対する演奏制御を実現する点にある．MIDI クロック同期や DMX512 のような専用規格を用いるのではなく，授業やハッカソンで扱いやすい Arduino と Processing を組み合わせることで，低コストで再現しやすい構成を実現する．また，Master がすべての演奏処理を集中管理するのではなく，各 Slave が担当楽器の処理を分担することで，処理負荷を分散できる．

さらに，本システムでは，演奏前に BPM，オクターブ，演奏順番を設定できるだけでなく，演奏中にも BPM を変更できる．この点は，単なる自動演奏装置ではなく，使用者が演奏の進行に関与できるインタラクティブな合奏システムとしての独自性である．また，BPM に連動したクマのアニメーションを Processing 上に表示することで，音楽と視覚情報を連動させたエンターテインメント性を付加している．これにより，楽器演奏の経験がない使用者でも，GUI 操作を通じて合奏を制御する楽しさを体験できる．

以上のように，本システムは，既存の音楽同期技術や1対多制御技術の考え方を参考にしながら，Arduino 複数台構成に適した低コストで簡易な同期方式を採用している．I2C 通信には通信距離やノイズ耐性に関する制約があるものの，小規模な教育用・体験用システムとしては，配線の簡易性，実装のしやすさ，処理分担の明確さという利点がある．そのため，本プロジェクトの目的である「誰でも指揮者になれる」合奏体験を実現するための通信方式として，I2C は妥当な選択であると考えられる．
